% Options for packages loaded elsewhere
\PassOptionsToPackage{unicode}{hyperref}
\PassOptionsToPackage{hyphens}{url}
%
\documentclass[
  ignorenonframetext,
]{beamer}
\usepackage{pgfpages}
\setbeamertemplate{caption}[numbered]
\setbeamertemplate{caption label separator}{: }
\setbeamercolor{caption name}{fg=normal text.fg}
\beamertemplatenavigationsymbolsempty
% Prevent slide breaks in the middle of a paragraph
\widowpenalties 1 10000
\raggedbottom
\setbeamertemplate{part page}{
  \centering
  \begin{beamercolorbox}[sep=16pt,center]{part title}
    \usebeamerfont{part title}\insertpart\par
  \end{beamercolorbox}
}
\setbeamertemplate{section page}{
  \centering
  \begin{beamercolorbox}[sep=12pt,center]{part title}
    \usebeamerfont{section title}\insertsection\par
  \end{beamercolorbox}
}
\setbeamertemplate{subsection page}{
  \centering
  \begin{beamercolorbox}[sep=8pt,center]{part title}
    \usebeamerfont{subsection title}\insertsubsection\par
  \end{beamercolorbox}
}
\AtBeginPart{
  \frame{\partpage}
}
\AtBeginSection{
  \ifbibliography
  \else
    \frame{\sectionpage}
  \fi
}
\AtBeginSubsection{
  \frame{\subsectionpage}
}
\usepackage{amsmath,amssymb}
\usepackage{lmodern}
\usepackage{iftex}
\ifPDFTeX
  \usepackage[T1]{fontenc}
  \usepackage[utf8]{inputenc}
  \usepackage{textcomp} % provide euro and other symbols
\else % if luatex or xetex
  \usepackage{unicode-math}
  \defaultfontfeatures{Scale=MatchLowercase}
  \defaultfontfeatures[\rmfamily]{Ligatures=TeX,Scale=1}
\fi
\usetheme[]{Montpellier}
% Use upquote if available, for straight quotes in verbatim environments
\IfFileExists{upquote.sty}{\usepackage{upquote}}{}
\IfFileExists{microtype.sty}{% use microtype if available
  \usepackage[]{microtype}
  \UseMicrotypeSet[protrusion]{basicmath} % disable protrusion for tt fonts
}{}
\makeatletter
\@ifundefined{KOMAClassName}{% if non-KOMA class
  \IfFileExists{parskip.sty}{%
    \usepackage{parskip}
  }{% else
    \setlength{\parindent}{0pt}
    \setlength{\parskip}{6pt plus 2pt minus 1pt}}
}{% if KOMA class
  \KOMAoptions{parskip=half}}
\makeatother
\usepackage{xcolor}
\IfFileExists{xurl.sty}{\usepackage{xurl}}{} % add URL line breaks if available
\IfFileExists{bookmark.sty}{\usepackage{bookmark}}{\usepackage{hyperref}}
\hypersetup{
  hidelinks,
  pdfcreator={LaTeX via pandoc}}
\urlstyle{same} % disable monospaced font for URLs
\newif\ifbibliography
\setlength{\emergencystretch}{3em} % prevent overfull lines
\providecommand{\tightlist}{%
  \setlength{\itemsep}{0pt}\setlength{\parskip}{0pt}}
\setcounter{secnumdepth}{-\maxdimen} % remove section numbering
%\documentclass[unicode,10pt]{beamer}% 'unicode'が必要
\usepackage{luatexja}% 日本語を使う場合は必須
%\usepackage[japanese]{babel}	
%\usefonttheme[onlymath]{serif}
%\usepackage[T1]{fontenc} %これを使うと、ギリシャ文字が出ない
%\usepackage{textcomp}
%\usepackage[scale=1.0]{tgheros} % Sans serif font
%\usepackage[scaled]{beramono} % Monospace font
%\usepackage{luatexja-otf}
%\renewcommand{\kanjifamilydefault}{\gtdefault}
%\usepackage[haranoaji]{luatexja-preset}

% スライド番号の表示 %%%%%%%%%%%%%%%%%%%
\makeatletter
\setbeamertemplate{footline}{%
  \leavevmode%
  \hbox{%
  \begin{beamercolorbox}[wd=.5\paperwidth,ht=2.25ex,dp=1ex,right]{date in head/foot}%
    \insertframenumber{} \hspace*{2ex} % new version without total frames
  \end{beamercolorbox}}%
  \vskip0pt%
}
\makeatother
% スライド番号の表示 %%%%%%%%%%%%%%%%%%%
\ifLuaTeX
  \usepackage{selnolig}  % disable illegal ligatures
\fi

\title{第8章: 次の一歩}
\subtitle{今井耕介 著\\
『社会科学のためのデータ分析入門 (QSS)』}
\author{}
\date{\vspace{-2.5em}2026-03-09}

\begin{document}
\frame{\titlepage}

\hypertarget{ux518dux73feux53efux80fdux306aux7814ux7a76}{%
\section{8.1
再現可能な研究}\label{ux518dux73feux53efux80fdux306aux7814ux7a76}}

\begin{frame}{再現可能性 (Reproducibility) の重要性}
\protect\hypertarget{ux518dux73feux53efux80fdux6027-reproducibility-ux306eux91cdux8981ux6027}{}
\begin{itemize}
\tightlist
\item
  \textbf{再現可能な研究}:
  他の研究者が同じデータとコードを用いて、同じ結果を得られること。
\item
  透明性の確保と、科学的知見の蓄積のために不可欠。
\item
  ツール：

  \begin{itemize}
  \tightlist
  \item
    \textbf{R Markdown}: コード、結果、文章を一つのドキュメントに統合。
  \item
    \textbf{Git / GitHub}: バージョン管理とコードの共有。
  \end{itemize}
\end{itemize}
\end{frame}

\hypertarget{ux7814ux7a76ux8ad6ux6587ux306eux57f7ux7b46}{%
\section{8.2
研究論文の執筆}\label{ux7814ux7a76ux8ad6ux6587ux306eux57f7ux7b46}}

\begin{frame}{論文の構成}
\protect\hypertarget{ux8ad6ux6587ux306eux69cbux6210}{}
\begin{enumerate}
\tightlist
\item
  \textbf{イントロダクション}: 研究の問いと背景。
\item
  \textbf{データと測定}: どのようなデータを用い、変数をどう定義したか。
\item
  \textbf{データ分析}: 因果推論や予測の手法。
\item
  \textbf{結果}: 図表を用いた効果的な提示。
\item
  \textbf{結論}: 知見のまとめと限界。
\end{enumerate}
\end{frame}

\begin{frame}{結果の提示}
\protect\hypertarget{ux7d50ux679cux306eux63d0ux793a}{}
\begin{itemize}
\tightlist
\item
  表だけでなく、\textbf{図 (Graphs)} を効果的に用いる。
\item
  推定値だけでなく、その\textbf{不確実性 (Standard Errors, Confidence
  Intervals)} を必ず示す。
\end{itemize}
\end{frame}

\hypertarget{ux7814ux7a76ux502bux7406}{%
\section{8.3 研究倫理}\label{ux7814ux7a76ux502bux7406}}

\begin{frame}{データ分析における倫理}
\protect\hypertarget{ux30c7ux30fcux30bfux5206ux6790ux306bux304aux3051ux308bux502bux7406}{}
\begin{itemize}
\tightlist
\item
  \textbf{プライバシーの保護}: 個人情報の匿名化。
\item
  \textbf{研究の公正性}: 都合の悪い結果を隠さない（p-hacking の回避）。
\item
  \textbf{透明性}: データとコードの公開を推奨。
\end{itemize}
\end{frame}

\hypertarget{ux3055ux3089ux306aux308bux5b66ux7fd2ux306eux305fux3081ux306b}{%
\section{8.4
さらなる学習のために}\label{ux3055ux3089ux306aux308bux5b66ux7fd2ux306eux305fux3081ux306b}}

\begin{frame}{今後の学習の方向性}
\protect\hypertarget{ux4ecaux5f8cux306eux5b66ux7fd2ux306eux65b9ux5411ux6027}{}
\begin{enumerate}
\tightlist
\item
  \textbf{高度な統計学・計量経済学}:

  \begin{itemize}
  \tightlist
  \item
    固定効果モデル、パネルデータ分析、操作変数法など。
  \end{itemize}
\item
  \textbf{機械学習 (Machine Learning)}:

  \begin{itemize}
  \tightlist
  \item
    大規模データからの予測やパターンの発見。
  \end{itemize}
\item
  \textbf{プログラミングスキル}:

  \begin{itemize}
  \tightlist
  \item
    スケールアップのための効率的なコーディング。
  \end{itemize}
\end{enumerate}
\end{frame}

\hypertarget{ux307eux3068ux3081}{%
\section{8.5 まとめ}\label{ux307eux3068ux3081}}

\begin{frame}{本書のまとめ}
\protect\hypertarget{ux672cux66f8ux306eux307eux3068ux3081}{}
\begin{itemize}
\tightlist
\item
  データを R で操作し、現実世界の問いに答える基礎を学んだ。
\item
  \textbf{測定、予測、説明 (因果推論)}
  という3つの大きな目標を意識すること。
\item
  データ分析は、社会を理解し、より良い政策決定を行うための強力な道具である。
\item
  常に批判的思考を持ち、不確実性を忘れずに分析を続けること。
\end{itemize}
\end{frame}

\end{document}
